\section{Methodology}\label{sec:method}

\subsection{The overnight-stay networks}

The Ministry of Transport publishes, for each night, an estimate of how many residents of each zone slept in each other zone \citep{mitms2024methodology}. According to the ministry's methodology, the estimate comes from the records of a mobile-phone operator: each phone is given a home zone, where it usually spends the night, and a zone for each night; the counts are then scaled up to the resident population. I use the 2,435 zones of mainland Spain, so that the network covers one connected territory, and the six months of 2024 in the data used here: February, March, August, October, November and December. A file dated $d$ holds the night that ends on the morning of $d$; I always name a night by both dates, as in ``14--15 August''.

For each night I build a directed network. A link $i\to j$ exists when at least 10 residents of $i$ sleep in $j$, and carries the estimated number of people as its weight. Stays inside one's own zone are not links. Below 10 people the estimates are noisy and the many small links would make every network almost complete; with the threshold, the networks keep 69--83\,\% of the people who sleep away from home and 0.9--3.4\,\% of all possible undirected links, so they stay sparse. Because the threshold keeps more links on busy nights, every structural comparison is repeated on a second version in which each night keeps the same number of directed links as the ordinary night, 35,812, chosen by weight (Table~\ref{tab:basic_full}).

Three nights are the main networks: an ordinary weekday (13--14 February), the summer holiday (14--15 August) and All Saints (31 October--1 November). Easter (28--29 March) and Christmas Eve (24--25 December) are extra checks. Any one night could be unusual, so each result is repeated on its \emph{sister nights}: the 13 February nights ending on a Tuesday, Wednesday or Thursday, the 9 nights ending on 10--18 August and the 3 nights ending on 1--3 November. For each headline result I report in how many pairs of one holiday sister and one February sister it keeps its sign (117 pairs for summer, 39 for All Saints), and the main nights are repeated in 2022 and, except All Saints, in 2023.

Table~\ref{tab:basic} and Figure~\ref{fig:basic} (in the appendix) give the basic properties of the three networks, computed on their undirected version (zones $i$ and $j$ are linked if either sends people to the other). All three are small worlds: in the giant component two zones are on average about two links apart (2.00 in summer), close to Erd\H{o}s--R\'enyi graphs simulated with the same numbers of nodes and links (2.02; Table~\ref{tab:basic_full}), while their clustering is 0.55--0.59 against 0.01--0.03 in those graphs. Their degree distributions, shown as complementary cumulative distributions (CCDF) in Figure~\ref{fig:basic}a, are heavy-tailed, but the evidence does not favour a power law: a likelihood-ratio test on the in- and out-degree tails \citep{clauset2009power} cannot tell a power law from a log-normal on the ordinary and All Saints nights and prefers the log-normal in summer. All three are disassortative: the mean degree of a zone's neighbours, $k_{nn}(k)$, falls with its own degree $k$, so well-connected zones tend to link to zones of lower degree.

\begin{table}[t]
\centering
\small
\setlength{\tabcolsep}{3pt}
\caption{Basic properties of the overnight-stay networks}
\label{tab:basic}
\resizebox{\linewidth}{!}{\begin{tabular}{lrrr}
\toprule
 & \shortstack[r]{Ordinary weekday\\13–14 Feb} & \shortstack[r]{Summer holiday\\14–15 Aug} & \shortstack[r]{All Saints\\31 Oct–1 Nov} \\
\midrule
Nodes N in the giant component & 2,432 & 2,435 & 2,434 \\
Links L & 26,465 & 101,996 & 48,698 \\
Density (\%) & 0.89 & 3.44 & 1.64 \\
Average degree $\langle$k$\rangle$ & 21.7 & 83.8 & 40.0 \\
Maximum degree $k_{\max}$ / share of all zones & 1,499 / 62 \% & 2,233 / 92 \% & 2,008 / 82 \% \\
Median out- / in-degree & 7 / 7 & 14 / 27 & 10 / 16 \\
Reciprocity (rewired) & 0.52 (0.13) & 0.24 (0.09) & 0.48 (0.15) \\
Assortativity r (rewired) & $-$0.131 ($-$0.140) & $-$0.214 ($-$0.217) & $-$0.155 ($-$0.161) \\
Rich club $\phi$, top 50 (rewired) & 0.80 (0.66) & 0.96 (0.92) & 0.91 (0.81) \\
Rich-club clique $n_{\mathrm{clique}}$ & 17 & 17$^\dagger$ & 22 \\
Average path length $\ell$ & 2.44 & 2.00 & 2.16 \\
Expected in a random graph $\ell^*$ & 2.53 & 1.76 & 2.11 \\
Clustering $\langle$C$\rangle$ (rewired / random) & 0.554 (0.247 / 0.009) & 0.546$^\dagger$ (0.387 / 0.034) & 0.593 (0.298 / 0.016) \\
\bottomrule
\end{tabular}}
\par\smallskip
\begin{minipage}{\linewidth}\footnotesize
Undirected networks of 2,435 zones, one per night of 2024; median degrees and reciprocity use the directed links. In brackets: mean of 50 degree-preserving rewirings (for clustering, also 50 Erdős–Rényi graphs). $^\dagger$ The direction of the change is not the same on every pair of sister nights. Definitions and the full version: Table A1.
\end{minipage}
\end{table}



\subsection{The roots network}

Where residents were born comes from the 10\,\% sample of the 2024 census \citep{ine2025censo}, which gives, for each person, the municipality of residence (if it has more than 10,000 inhabitants), the province of birth, the size of the municipality of birth, and the year of arrival in the current province (not the city). I keep the 675 municipalities over 10,000 inhabitants that form a zone of their own and lie in one of the 44 provinces with at least two of them, so that each can be compared with a neighbour. For each municipality $c$ and each other mainland province $p$, the roots network counts the residents of $c$ born in $p$, in total and split into groups: by arrival (before 1981, which I use as a marker of the exodus generation, and 2016--2023) and by birthplace (a village of 2,000 inhabitants or fewer, or a town of more than 10,000). Each group share is the group's count of residents of $c$ born in $p$ over all sampled residents of $c$. I call these 675 municipalities \emph{cities}, although many are small towns, and use them for the group results. The overall pull of the roots, without groups, uses the full published census table of residents by province of birth, which covers 137 large mainland cities, each a zone of its own; 112 of them lie in the 22 provinces with at least two, which identify the comparison.

\subsection{Methods}

\paragraph{Structure against degree-preserving null models (T1, T2).} The holiday networks are much denser, and a denser network looks different on nearly every measure just because it has more links. So I compare every property with a degree-preserving null model, which I use throughout the report. I take each network and repeatedly swap the ends of two random links (10 attempted swaps per link), so that every zone keeps its number of links but who is linked to whom becomes random. I report the real value next to the mean of 50 of these: if the two are close, the degrees alone explain the property. For degree mixing I use the rewiring of \citet{zhou2007structural}: pushing it towards the most assortative and the most disassortative network it can reach, while keeping every degree and all zones with links in one connected component, shows the range of $r$ that the degrees allow and where the real network sits in it. The rich club is measured as $\phi(r)$ against the rank $r$ of the best-connected zones \citep{zhou2004richclub}, raw and against the rewired networks. Second-order mixing $R_{\mathrm{avg}}$ \citep{zhou2017second} goes one step further and asks whether the neighbours of two linked zones are alike: for each end of a link I average the excess degrees ($k-1$) of its other neighbours, using zero if there are none, and correlate the two ends ($R_{\max}$ uses the largest excess degree instead of the average). Its null keeps the degrees and also the rich-club coefficient $\phi(k)$ at every threshold (link-rewiring algorithm No.~2 of the course). The appendix gives the bitmap of the sorted adjacency matrix (BOSAM) of each night \citep{guo2009fingerprint}.

\paragraph{The pull of the roots (T3).} This was the part that took me longest. Correlating the two networks, as \citet{li2016spring} do, gives a high value, but both fall with distance (Section~\ref{sec:literature}): people tend to be born near where they live, and also to travel near where they live.

So I turned it into a comparison of like with like: two cities of the same province, sending people to the same destination province. Anything that belongs to the pair of provinces (how far apart they are, whether the destination has a coast) is the same for both cities and cancels out; on top of that I model the distance from each city. The outcome adds up all of the city's residents who sleep anywhere in $p$, including flows below the 10-person threshold. These are counts with many zeros, so I used Poisson pseudo-maximum likelihood (PPML), a regression for counts that keeps the zeros. For city $c$ in province $o(c)$, destination province $p\neq o(c)$ and night $t$, the number of residents of $c$ who sleep in $p$, $y_{cpt}$, is modelled as
\begin{equation}
\mathrm{E}[y_{cpt}] = \exp\big(\beta_t \log s_{cp} + \alpha_{ct} + \gamma_{o(c)pt} + f_t(d_{cp}) + \lambda_{pt}\,\log I_c\big),
\label{eq:ppml}
\end{equation}
where $s_{cp}$ is the share of the city's residents born in $p$, $\alpha_{ct}$ is a separate baseline for each city on each night (how much its residents travel at all), $\gamma_{o(c)pt}$ a separate baseline for each ordered pair of origin and destination provinces on each night, $f_t$ is a quadratic in the log distance from the city to the centre of $p$, and the last term lets the income per person of the city, $I_c$ \citep{ine2026adrh}, affect each destination differently. The number I care about is $\beta_t$, the \emph{pull of the roots}: with everything else fixed, a city with twice the share of residents born in $p$ is expected to send $2^{\beta_t}$ times as many people there. As a link-by-link check closer to \citet{li2016spring}, I fit to each network (the roots and each night's stays) a model that keeps each city's total, each province's total and the total in each distance band, and correlate the logged ratios of observed to expected links of the two; the $p$-value comes from permuting cities within their province (QAP, 2,000 permutations).

To see whether the pull changes on holiday, I fit two nights together, each with its own baselines and distance terms, and estimate the change $\delta=\beta_{\text{holiday}}-\beta_{\text{ordinary}}$ directly. To see \emph{who} carries it, I replace $\log s_{cp}$ by the shares of two groups at a time: arrived before 1981 and arrived in 2016--2023 in one model, born in a village and born in a town in another. One practical problem was that in the 10\,\% sample many city--province cells are empty for a group, so their log is undefined. I add a half to every group count before dividing by the city's residents; Table~\ref{tab:zero} shows that five other fixes give the same signs, though not always the same sizes. Each model is also fitted separately for stays in inland and in coastal zones of $p$.

With only 22 provinces for the 112 large cities and 44 for the 675, the usual standard errors are too optimistic, so every $p$-value and interval comes from a wild cluster bootstrap by province of origin with 9,999 draws \citep{kline2012score}, which lets cities of the same province have related errors. The main baseline is the February night; to check that a winter baseline is not behind the results, I added a second one, the mean of the 15 Tuesday--Thursday nights of October.

\paragraph{The emptied villages (T4).} For each zone I measure how much population it lost: one minus the ratio of its 2021 population to the largest population it had between 1900 and 1960, from the historical census series \citep{ine2026intercensal}. Zones are split into five groups of about 475 by this loss. A zone is a \emph{net receiver} on a night when more people from elsewhere sleep in it than residents of the zone sleep elsewhere, using all flows, not only links above the threshold; its \emph{net balance} is (in $-$ out)/(in $+$ out). I compare the share of net receivers in each group across nights, and on each night I regress being a net receiver on the percentile rank of loss (0 for the least emptied zone, 1 for the most), with zone population, number of municipalities, area, coast, distance to the coast and to the nearest city of 100,000 people, and province as controls. To see whether second homes are part of the story, I also regress the log ratio of visitors on a holiday night to February on the same rank, with and without the share of non-primary dwellings \citep{ine2023viviendas}. Three zones with gaps in the data and 56 without population history are excluded, leaving 2,376.

